\section{Heat smoothing and entry from spatial initial metrics}
\label{sec:spatial-entry}

Fix a smooth base point $p$ and an actual HE representative
$w$ from Proposition~\ref{prop:exact-balancing-coordinates}.
Throughout the flow argument these parameters remain fixed.
Write $r=|w|$, $D=D_{\mathcal A_p(w)}=D_0+a$ and
$L=D^*D$, $L_0=D_0^*D_0$. All section variables are
Hermitian and trace free for the fixed transported metric $h_0$.
Let $\Pi$ be the split projection onto $\mathfrak h$ and
$Q=I-\Pi$. On each smooth fibre put
\begin{equation}\label{eq:flow-banach-spaces}
 Y_p=C^0(\Herm_0E^{\rm sp}),\qquad
 X_p=C^1(\Herm_0E^{\rm sp}),\qquad
 \|u\|_X=\|u\|_\infty+\|D_0u\|_\infty.
\end{equation}
We use these fibrewise spaces with uniform operator bounds;
no fixed Banach bundle through the singular fibre is required.

\begin{theorem}\label{thm:spatial-entry}
There are $M,\sigma,C,\delta,r_0>0$, independent of the
base scales and of all positive weak-rate ratios, such that
\begin{equation}\label{eq:split-heat-smoothing}
 \begin{split}
 \|e^{-tL_0}Q\|_{X_p\to X_p}&\le Me^{-\sigma t},\\
 \|e^{-tL_0}Q\|_{Y_p\to X_p}
                    &\le M(1+t^{-1/2})e^{-\sigma t}\quad(t>0).
 \end{split}
\end{equation}
In the metric coordinate $h=h_0e^u$, the actual HYM equation is
\begin{equation}\label{eq:exact-log-flow}
                  \dot u+L_0u=N_w(u).
\end{equation}
Here $N_w:X_p\to Y_p$ depends only on the first spatial
jet, and is smooth with uniform derivatives of every fixed
finite order on a common small ball. For parallel $S$ and
$v\in QX_p$ with $|S|,\|v\|_X\le\delta_1$, it satisfies
\begin{equation}\label{eq:flow-nonlinearity-bounds}
 \begin{split}
 \|N_w(S)\|_\infty&\le Cr^2|S|,\\
 \|N_w(S+v)-N_w(S)\|_\infty&\le C(r+\delta_1)\|v\|_X.
 \end{split}
\end{equation}
Every fixed positive-order derivative of $N_w(S)$ in $S$
is $O(r^2)$ on this ball.

For every smooth $u_0$ with $\|u_0\|_X<\delta$ and
$r<r_0$, the determinant-normalized actual flow exists
for all $t\ge0$, stays in a common small $C^1$ ball, and obeys
\begin{equation}\label{eq:normal-entry}
 \begin{split}
 \|u(t)\|_\infty&\le\|u_0\|_\infty,
       \qquad |\Pi u(t)|\le\|u_0\|_\infty,\\
 \|Qu(t)\|_X&\le Ce^{-\sigma t/2}\|Qu_0\|_X
                         +Cr^2\|u_0\|_\infty.
 \end{split}
\end{equation}
For $0<r<\min(r_0,1)$, the time
$T_r=4\sigma^{-1}\log(1/r)$ gives $\|Qu(t)\|_X\le Cr^2\delta$
for all $t\ge T_r$. Every positive weak eigenvalue satisfies
$T_r\lambda_j\to0$ uniformly as $r\to0$.
\end{theorem}

\subsection{Uniform first derivatives of the split heat flow}

Let $S_t=e^{-tL_0}$ and let $H_t$ denote scalar heat flow.
Kato domination gives $|S_tf|\le H_t|f|$, so
$\|S_tf\|_\infty\le\|f\|_\infty$.
The scalar estimates from Section~\ref{sec:projected-he} and
Bochner's formula with $\operatorname{Ric}(g_p)=0$ give
\begin{equation}\label{eq:scalar-two-gradient-bounds}
 \|\dd H_t\phi\|_\infty\le Ct^{-1/2}\|\phi\|_\infty,
 \qquad
 \|\dd H_t\phi\|_\infty\le C\|\dd\phi\|_\infty.
\end{equation}
Use the fixed finite cover and unitary line frames from
the proof of Lemma~\ref{lem:graph-inverse}. Their connection
coefficients are bounded and divergence free. The partition
functions $\chi_j$ have uniformly bounded gradients and
Laplacians; the derivatives of the core cutoffs are supported
in fixed exterior annuli. If $f_j(t)$ is a scalar coefficient
of $S_tf$ in one frame, extend $z_j=\chi_jf_j$ by zero
to the whole fibre. The local formula
\eqref{eq:local-rough-operator} gives
\[
 (\partial_t+\Delta)z_j=F_j,
 \qquad \|F_j(t)\|_\infty
       \le C\{\|D_0S_tf\|_\infty+\|f\|_\infty\}.
\]
Conversely,
$\|D_0S_tf\|_\infty\le C\sum_j\|\dd z_j\|_\infty+C\|f\|_\infty$.
No derivative of the full curvature enters these formulas.

Scalar Duhamel and \eqref{eq:scalar-two-gradient-bounds}
therefore give, for $G(t)=\|D_0S_tf\|_\infty$ and
$F=\|f\|_\infty$,
\[
 G(t)\le C(t^{-1/2}+1)F+
      C\int_0^t(t-s)^{-1/2}\{G(s)+F\}\,\dd s
                       \quad(0<t\le1).
\]
For smooth $f$, set $Z_T=\sup_{0<t\le T}\sqrt t\,G(t)$.
The identity
$\int_0^t(t-s)^{-1/2}s^{-1/2}\,\dd s=\pi$
implies $Z_T\le CF+C\sqrt T Z_T$. One common small
$T$ absorbs the last term. The semigroup law extends the
bound to fixed bounded time intervals. For $f\in C^1$,
use the second scalar estimate on the initial term instead.
Taking $\sup_{0<t\le T}G(t)$ and using
$\int_0^t(t-s)^{-1/2}\,\dd s=2\sqrt t$ yields
\[
 \|S_tf\|_X\le C\|f\|_X\quad(0\le t\le1),\qquad
 \|S_tf\|_X\le C(1+t^{-1/2})\|f\|_\infty\quad(0<t\le1).
\]
Density extends these estimates. On each fixed smooth fibre,
the smooth heat solutions converge to their initial data in
$C^1$; density and the first estimate prove strong continuity
on $X_p$. The same argument works on $Y_p$.

The averaging projection $\Pi$ is uniformly bounded on both
spaces and commutes with $S_t$. The split $L^2$ gap, fixed
volume, and the fixed-time $L^2$-to-$L^\infty$ bound used in
Lemma~\ref{lem:graph-inverse} give
$\|S_tQ\|_{Y_p\to Y_p}\le Ce^{-ct}$ for $t\ge1$.
Composing with a fixed positive-time $Y_p$-to-$X_p$ map
proves \eqref{eq:split-heat-smoothing}. This argument supplies
the first-derivative smoothing needed below without a
uniform higher-order elliptic estimate.

\subsection{The contracted connection derivative}

The connection expansion already gives $\|a\|_\infty\le Cr$.
We also need its divergence:
\begin{equation}\label{eq:connection-divergence}
                         \|D_0^*a\|_\infty\le Cr^2.
\end{equation}
Indeed put $U=e^{s_p(w)/2}$ and
$\gamma=Ub_p(w)U^{-1}-(\dbar_0U)U^{-1}$, so $a=\gamma-\gamma^*$.
Section~\ref{sec:projected-he} gives
\[
 \|b_p(w)\|_\infty\le Cr,\quad
 \|\dbar_0^*b_p(w)\|_\infty\le Cr^2,
 \quad \|D_0U\|_\infty+\|P_pU\|_\infty\le Cr^2.
\]
The same first-derivative estimate holds for $U^{-1}$.
The divergence product rule gives
\[
 \begin{split}
 \dbar_0^*(UbU^{-1})
 &=U(\dbar_0^*b)U^{-1}
                  +O_\infty(|D_0U||b|+|b||D_0U^{-1}|),\\
 \dbar_0^*((\dbar_0U)U^{-1})
 &=(P_pU)U^{-1}+O_\infty(|D_0U||D_0U^{-1}|).
 \end{split}
\]
Thus $\|\dbar_0^*\gamma\|_\infty\le Cr^2$.
Real unitary divergence commutes with adjoints, proving
\eqref{eq:connection-divergence}.
For $A_j=\operatorname{ad}_{a(e_j)}$ in a real orthonormal
frame, expansion of the connection Laplacian gives
\begin{equation}\label{eq:rough-operator-perturbation}
 (L-L_0)u=-2\sum_j A_jD_{0,e_j}u+[D_0^*a,u]-\sum_jA_j^2u.
\end{equation}
Its first-order coefficient is $O(r)$ and its zeroth-order
coefficient is $O(r^2)$.

\subsection{The exact equation in logarithmic coordinates}

Define the pointwise invertible linear map
\[
 J_u(V)=\int_0^1e^{-\theta u}Ve^{\theta u}\,\dd\theta.
\]
It satisfies $e^{-u}(\dd\exp)_u(V)=J_u(V)$.
Diagonalization of $u$ shows that its scalar multipliers
are $(1-e^{-a})/a$, with value one at $a=0$.
Its inverse and all fixed-order derivatives are bounded
on a fixed small pointwise ball. Since the reference
connection $D$ is HE, the exact Chern formula is
\[
 \mathrm i\Lambda F_{h_0e^u}
 =\mathrm i\Lambda\dbar_D(J_u(D'_Du))
 =J_u(\mathrm i\Lambda\dbar_DD'_Du)
       +\mathrm i\Lambda(dJ)_u[\dbar_Du]\wedge D'_Du.
\]
The section Kähler identity gives
$\mathrm i\Lambda\dbar_DD'_Du=Lu/2$.
Consequently $h^{-1}\dot h=-2\mathrm i\Lambda F_h$ becomes
\begin{equation}\label{eq:first-jet-nonlinearity}
 \begin{split}
 \dot u+Lu&=\mathcal Q(u,Du),\\
 \mathcal Q(u,Du)&=-2J_u^{-1}\mathrm i\Lambda
                       (dJ)_u[\dbar_Du]\wedge D'_Du.
 \end{split}
\end{equation}
Here $dJ$ differentiates its finite matrix coefficient $u$.
The nonlinear term is quadratic in the first spatial jet;
all second derivatives have cancelled. Thus
\begin{equation}\label{eq:Nw-definition}
 N_w(u)=-(L-L_0)u+\mathcal Q(u,D_0u+[a,u])
\end{equation}
has the asserted uniform smoothness from $X_p$ to $Y_p$.
This follows from bounded matrix derivatives and the
continuous pointwise product $C^0\times C^0\to C^0$.
The equation preserves Hermitian variables by its metric
origin. It preserves trace zero because the determinant
metric solves scalar heat and initially equals the fixed
HE determinant.

For parallel $S$, the first-order term in
\eqref{eq:rough-operator-perturbation} vanishes, and
$DS=[a,S]$. Hence $N_w(S)=O(r^2|S|)+O(r^2|S|^2)$,
with the same factor $r^2$ in every fixed finite derivative.
Taking the difference at $S+v$ proves
\eqref{eq:flow-nonlinearity-bounds}.

\subsection{Confinement and normal decay}

On each fixed compact smooth fibre, global smooth HYM
metric flow from a smooth positive metric follows from
\cite[Theorem 6.8.25]{Kobayashi}. We use this only for existence;
its higher constants need not be uniform. The squared
pointwise distance from $h_0$ to $h_0e^u$ in the space
of positive Hermitian metrics is $\tr u^2$.
The two-solution distance inequality
\cite[Lemma 6.4.5]{Kobayashi}, with time changed by $\tau=2t$
to our convention, gives
\begin{equation}\label{eq:metric-distance-heat}
          (\partial_t+\Delta)\tfrac12\tr u^2\le0.
\end{equation}
The reference HE metric is stationary. The maximum principle
therefore proves the first bound in \eqref{eq:normal-entry}.
Block averaging and Jensen's inequality give the bound
for $S(t)=\Pi u(t)$.

Let $v=Qu$ and $\varepsilon=\|u_0\|_\infty$.
As long as $\|v\|_X\le\delta_1$, Duhamel gives
\begin{equation}\label{eq:normal-volterra}
 \|v(t)\|_X\le Me^{-\sigma t}\|v(0)\|_X+
  \int_0^t k(t-s)\{a_0\|v(s)\|_X+b_0\}\,\dd s,
\end{equation}
where $k(t)=M'(1+t^{-1/2})e^{-\sigma t}$,
$a_0=C(r+\delta_1)$, and $b_0=Cr^2\varepsilon$.
Choose $\delta_1,r_0$ so that
$a_0\int_0^\infty k(t)e^{\sigma t/2}\,\dd t\le1/2$.
Iteration of this nonnegative convolution inequality gives
\[
 \|v(t)\|_X\le2Me^{-\sigma t/2}\|v(0)\|_X
                              +2b_0\int_0^\infty k(s)\,\dd s.
\]
Equivalently the right side is a supersolution: its
exponential term is controlled by the weighted kernel
integral and its constant term by the unweighted integral.
Uniform boundedness of $\Pi,Q$ lets us choose one initial
$\delta$ making this bound smaller than $\delta_1/2$.
A first-exit argument then proves it for all time.
Since $D_0S=0$, this also confines $u$ in the required
$C^1$ ball and completes \eqref{eq:normal-entry}.

At $T_r$ the exponential is at most $r^2$.
Theorem~\ref{thm:complete-bundle-spectrum} gives
$\lambda_j\le Cr^2$ for every positive weak rate, so
$T_r\lambda_j\le Cr^2\log(1/r)\to0$.
This proves Theorem~\ref{thm:spatial-entry}, including
arbitrarily unequal weak rates. At $r=0$ there are no
positive weak rates, and the normal part decays exponentially.

The basin is an open ball in the $C^1$ topology intersected
with smooth determinant-normalized metrics. It contains
spatially varying initial metrics. The persistent $O(r^2)$
normal term in \eqref{eq:normal-entry} does not by itself
give tracking on a long time interval. We next absorb
that term into an invariant finite-dimensional graph.
